\begin{thebibliography}{99}
\bibitem{Antoniano1985} J. L. Antoniano \& G. A. Uhlmann - “A functional calculus for a class of pseudodifferential operators with singular symbols” , in Pseudodifferential operators and applications (Notre Dame, Ind., 1984) , Proc. Sympos. Pure Math. , vol. 43 , American Mathematical Society, Providence, RI, 1985, p. 5-16
\bibitem{Baladi2018} V. Baladi - Dynamical zeta functions and dynamical determinants for hyperbolic maps , Ergeb. Math. Grenzgeb. (3) , vol. 68 , Springer, Cham, 2018
\bibitem{Bar2019} C. Bär \& A. Strohmaier - “An index theorem for Lorentzian manifolds with compact spacelike Cauchy boundary” , Amer. J. Math. 141 (2019) no. 5, p. 1421-1455
\bibitem{Baer2020} C. Bär \& A. Strohmaier - “Local index theory for Lorentzian manifolds” , 2020
\bibitem{Bischo2020} F. Bischoff, H. Bursztyn, H. Lima \& E. Meinrenken - “Deformation spaces and normal forms around transversals” , Compositio Math. 156 (2020) no. 4, p. 697-732
\bibitem{Bourgain2015} J. Bourgain, P. Shao, C. D. Sogge \& X. Yao - “On $L^p$-resolvent estimates and the density of eigenvalues for compact Riemannian manifolds” , Comm. Math. Phys. 333 (2015) no. 3, p. 1483-1527
\bibitem{Brunetti2000} R. Brunetti \& K. Fredenhagen - “Microlocal analysis and interacting quantum field theories: renormalization on physical backgrounds” , Comm. Math. Phys. 208 (2000) no. 3, p. 623-661
\bibitem{Bursztyn2019} H. Bursztyn, H. Lima \& E. Meinrenken - “Splitting theorems for Poisson and related structures” , J. reine angew. Math. (2019) no. 754, p. 281-312
\bibitem{Connes1988} A. Connes - “The action functional in non-commutative geometry” , Comm. Math. Phys. 117 (1988) no. 4, p. 673-683
\bibitem{Connes1995} A. Connes \& H. Moscovici - “The local index formula in noncommutative geometry” , Geom. Funct. Anal. 5 (1995) no. 2, p. 174-243
\bibitem{dangthesis} N. V. Dang - Renormalization of quantum field theory on curved space-times, a causal approach , Ph. D. Thesis, Université Paris Diderot (Paris VII), 2013
\bibitem{DangAHP} N. V. Dang - “The extension of distributions on manifolds, a microlocal approach” , Ann. Henri Poincaré 17 (2016) no. 4, p. 819-859
\bibitem{Dang2020} N. V. Dang \& M. Wrochna - “Complex powers of the wave operator and the spectral action on Lorentzian scattering spaces” , 2020
\bibitem{derezinski} J. Dereziński \& D. Siemssen - “Feynman propagators on static spacetimes” , Rev. Math. Phys. 30 (2018) no. 3, article ID 1850006, 23 pages
\bibitem{Derezinski2019} J. Dereziński \& D. Siemssen - “An evolution equation approach to linear quantum field theory” , 2019
\bibitem{DeWitt1975} B. S. DeWitt - “Quantum field theory in curved spacetime” , Phys. Rep. 19 (1975) no. 6, p. 295-357
\bibitem{Ferreira2014} D. Dos Santos Ferreira, C. E. Kenig \& M. Salo - “On $L^p$ resolvent estimates for Laplace–Beltrami operators on compact manifolds” , Forum Math. 26 (2014) no. 3, p. 815-849
\bibitem{Dowker1976} J. S. Dowker \& R. Critchley - “Effective Lagrangian and energy-momentum tensor in de Sitter space” , Phys. Rev. D 13 (1976) no. 12, p. 3224-3232
\bibitem{Dyatlov2016} S. Dyatlov \& M. Zworski - “Dynamical zeta functions for Anosov flows via microlocal analysis” , Ann. Sci. École Norm. Sup. (4) 49 (2016) no. 3, p. 543-577
\bibitem{Fulling1989} S. A. Fulling - Aspects of quantum field theory in curved space-time , London Math. Soc. Student Texts , vol. 17 , Cambridge University Press, Cambridge, 1989
\bibitem{GHV} J. Gell-Redman, N. Haber \& A. Vasy - “The Feynman propagator on perturbations of Minkowski space” , Comm. Math. Phys. 342 (2016) no. 1, p. 333-384
\bibitem{GWfeynman} C. Gérard \& M. Wrochna - “The massive Feynman propagator on asymptotically Minkowski spacetimes” , Amer. J. Math. 141 (2019) no. 6, p. 1501-1546
\bibitem{Gerard2019b} C. Gérard \& M. Wrochna - “The massive Feynman propagator on asymptotically Minkowski spacetimes II” , Internat. Math. Res. Notices (2020) no. 20, p. 6856-6870
\bibitem{Gracia-Bondia2014} J. M. Gracia-Bondía, H. Gutiérrez \& J. C. Várilly - “Improved Epstein–Glaser renormalization in x-space versus differential renormalization” , Nuclear Phys. B 886 (2014), p. 824-869
\bibitem{Greenleaf1990} A. Greenleaf \& G. Uhlmann - “Estimates for singular Radon transforms and pseudodifferential operators with singular symbols” , J. Funct. Anal. 89 (1990) no. 1, p. 202-232
\bibitem{Guillemin1985} V. Guillemin - “A new proof of Weyl’s formula on the asymptotic distribution of eigenvalues” , Adv. Math. 55 (1985) no. 2, p. 131-160
\bibitem{Guillemin1993} V. Guillemin - “Gauged Lagrangian distributions” , Adv. Math. 102 (1993) no. 2, p. 184-201
\bibitem{Guillemin1981} V. Guillemin \& G. Uhlmann - “Oscillatory integrals with singular symbols” , Duke Math. J. 48 (1981) no. 1, p. 251-267
\bibitem{Hack2012a} T.-P. Hack \& V. Moretti - “On the stress–energy tensor of quantum fields in curved spacetimes – comparison of different regularization schemes and symmetry of the Hadamard/Seeley–DeWitt coefficients” , J. Phys. A 45 (2012) no. 37, article ID 374019
\bibitem{Hartung2015} T. Hartung \& S. Scott - “A generalized Kontsevich-Vishik trace for Fourier integral operators and the Laurent expansion of $\zeta $-functions” , 2015
\bibitem{Hawking1977} S. W. Hawking - “Zeta function regularization of path integrals in curved spacetime” , Comm. Math. Phys. 55 (1977) no. 2, p. 133-148
\bibitem{Hollands2001} S. Hollands \& R. M. Wald - “Local Wick polynomials and time ordered products of quantum fields in curved spacetime” , Comm. Math. Phys. 223 (2001) no. 2, p. 289-326
\bibitem{H} L. Hörmander - The analysis of linear partial differential operators I. Distribution theory and Fourier analysis , Classics in Math. , Springer-Verlag, Berlin, 2003
\bibitem{HormanderIII} L. Hörmander - The analysis of linear partial differential operators III. Pseudo-differential operators , Classics in Math. , Springer, Berlin, 2007
\bibitem{10.2307/2162103} M. S. Joshi - “An intrinsic characterisation of polyhomogeneous Lagrangian distributions” , Proc. Amer. Math. Soc. 125 (1997) no. 5, p. 1537-1543
\bibitem{joshi} M. S. Joshi - “Complex powers of the wave operator” , Portugal. Math. 54 (1997) no. 3, p. 345-362
\bibitem{Joshi1998} M. S. Joshi - “A symbolic construction of the forward fundamental solution of the wave operator” , Comm. Partial Differential Equations 23 (1998) no. 7-8, p. 1349-1417
\bibitem{Kay1991} B. S. Kay \& R. M. Wald - “Theorems on the uniqueness and thermal properties of stationary, nonsingular, quasifree states on spacetimes with a bifurcate Killing horizon” , Phys. Rep. 207 (1991) no. 2, p. 49-136
\bibitem{Kontsevich1995} M. Kontsevich \& S. Vishik - “Geometry of determinants of elliptic operators” , in Functional analysis on the eve of the 21st century, Vol. 1 (New Brunswick, NJ, 1993) , Progress in Math. , vol. 131 , Birkhäuser Boston, Boston, MA, 1995, p. 173-197
\bibitem{Lesch} M. Lesch - “On the noncommutative residue for pseudodifferential operators with log-polyhomogeneous symbols” , Ann. Global Anal. Geom. 17 (1999), p. 151-187
\bibitem{Lesch2000} M. Lesch \& M. J. Pflaum - “Traces on algebras of parameter dependent pseudodifferential operators and the eta–invariant” , Trans. Amer. Math. Soc. 352 (2000) no. 11, p. 4911-4936
\bibitem{Lewandowski2020} M. Lewandowski - “Hadamard states for bosonic quantum field theory on globally hyperbolic spacetimes” , 2020
\bibitem{Maeda2005} Y. Maeda, D. Manchon \& S. Paycha - “Stokes’ formulae on classical symbol valued forms and applications” , 2005
\bibitem{Meinrenken2021} E. Meinrenken - “Euler-like vector fields, normal forms, and isotropic embeddings” , Indag. Math. (N.S.) 32 (2021) no. 1, p. 224-245
\bibitem{Melrose1979} R. B. Melrose \& G. A. Uhlmann - “Lagrangian intersection and the Cauchy problem” , Comm. Pure Appl. Math. 32 (1979) no. 4, p. 483-519
\bibitem{Meyer} Y. Meyer - Wavelets, vibrations and scalings , CRM Monograph Series , vol. 9 , American Mathematical Society, Providence, RI, 1998
\bibitem{Minakshisundaram1949} S. Minakshisundaram \& Å. Pleijel - “Some properties of the eigenfunctions of the Laplace-operator on Riemannian manifolds” , Canad. J. Math. 1 (1949) no. 3, p. 242-256
\bibitem{Moretti1999} V. Moretti - “Local $\zeta $-function techniques vs. point-splitting procedure: A few rigorous results” , Comm. Math. Phys. 201 (1999) no. 2, p. 327-363
\bibitem{morettilong} V. Moretti - “One-loop stress-tensor renormalization in curved background: The relation between $\zeta $-function and point-splitting approaches, and an improved point-splitting procedure” , J. Math. Phys. 40 (1999) no. 8, p. 3843-3875
\bibitem{nakamurataira} S. Nakamura \& K. Taira - “Essential self-adjointness of real principal type operators” , Ann. H. Lebesgue 4 (2021), p. 1035-1059
\bibitem{Nikolov2013} N. M. Nikolov, R. Stora \& I. Todorov - “Euclidean configuration space renormalization, residues and dilation anomaly” , in Lie theory and its applications in physics , Springer Proc. Math. Stat. , vol. 36 , Springer, Tokyo, 2013, p. 127-147
\bibitem{paycha2} S. Paycha - “The noncommutative residue and canonical trace in the light of Stokes’ and continuity properties” , 2007
\bibitem{paycha} S. Paycha - Regularised integrals, sums and traces. An analytic point of view , University Lecture Series , vol. 59 , American Mathematical Society, Providence, RI, 2012
\bibitem{Pollicott1986} M. Pollicott - “Meromorphic extensions of generalised zeta functions” , Invent. Math. 85 (1986) no. 1, p. 147-164
\bibitem{Radzikowski1996} M. J. Radzikowski - “Micro-local approach to the Hadamard condition in quantum field theory on curved space-time” , Comm. Math. Phys. 179 (1996) no. 3, p. 529-553
\bibitem{Rejzner2020} K. Rejzner - “Renormalization and periods in perturbative algebraic quantum field theory” , in Periods in quantum field theory and arithmetic , Springer Proc. Math. Stat. , vol. 314 , Springer, Cham, 2020, p. 345-376
\bibitem{Ruelle1986} D. Ruelle - “Resonances of chaotic dynamical systems” , Phys. Rev. Lett. 56 (1986) no. 5, p. 405-407
\bibitem{seeley} R. T. Seeley - “Complex powers of an elliptic operator” , in Singular Integrals (Chicago, Ill., 1966) , Proc. Symp. Pure Math. , vol. 10 , American Mathematical Society, Providence, RI, 1967, p. 288-307
\bibitem{Shen2021} D. Shen \& M. Wrochna - “An index theorem on asymptotically static spacetimes with compact Cauchy surface” , 2021
\bibitem{shubin} M. A. Shubin - Pseudodifferential operators and spectral theory , Springer-Verlag, Berlin, 2001
\bibitem{Sogge1988} C. D. Sogge - “Concerning the $L^p$ norm of spectral clusters for second-order elliptic operators on compact manifolds” , J. Funct. Anal. 77 (1988) no. 1, p. 123-138
\bibitem{soggeHangzhou} C. D. Sogge - Hangzhou lectures on eigenfunctions of the Laplacian , Annals of Math. Studies , vol. 188 , Princeton University Press, Princeton, NJ, 2014
\bibitem{Strohmaier2020b} A. Strohmaier \& S. Zelditch - “A Gutzwiller trace formula for stationary space-times” , Adv. Math. 376 (2020), article ID 107434, 53 pages
\bibitem{Strohmaier2020a} A. Strohmaier \& S. Zelditch - “Spectral asymptotics on stationary space-times” , Rev. Math. Phys. (2020), article ID X206000
\bibitem{Strohmaier2020} A. Strohmaier \& S. Zelditch - “Semi-classical mass asymptotics on stationary spacetimes” , Indag. Math. (N.S.) 32 (2021) no. 1, p. 323-363
\bibitem{Taira2020a} K. Taira - “Limiting absorption principle and equivalence of Feynman propagators on asymptotically Minkowski spacetimes” , Comm. Math. Phys. 388 (2021) no. 1, p. 625-655
\bibitem{Vasy2017b} A. Vasy - “On the positivity of propagator differences” , Ann. Henri Poincaré 18 (2017) no. 3, p. 983-1007
\bibitem{vasyessential} A. Vasy - “Essential self-adjointness of the wave operator and the limiting absorption principle on Lorentzian scattering spaces” , J. Spectral Theory 10 (2020) no. 2, p. 439-461
\bibitem{vasywrochna} A. Vasy \& M. Wrochna - “Quantum fields from global propagators on asymptotically Minkowski and extended de Sitter spacetimes” , Ann. Henri Poincaré 19 (2018) no. 5, p. 1529-1586
\bibitem{Wald1979} R. M. Wald - “On the Euclidean approach to quantum field theory in curved spacetime” , Comm. Math. Phys. 70 (1979) no. 3, p. 221-242
\bibitem{wodzicki} M. Wodzicki - “Local invariants of spectral asymmetry” , Invent. Math. 75 (1984) no. 1, p. 143-177
\bibitem{StevenZelditch2012} S. Zelditch - “Pluri-potential theory on Grauert tubes of real analytic Riemannian manifolds, I” , in Spectral geometry , Proc. Sympos. Pure Math. , vol. 84 , American Mathematical Society, Providence, RI, 2012, p. 299-339
\bibitem{Zelditch2017} S. Zelditch - Eigenfunctions of the Laplacian of Riemannian manifolds , CBMS Regional Conference Series in Math. , vol. 125 , American Mathematical Society, Providence, RI, 2017
\end{thebibliography}
